% Owner: A
We start from the bipartite graph that links each recipe to its ingredients
and project it onto the ingredients, as in the course~\cite{quercini2026}. In
the course projection, two ingredients are linked with a weight equal to
their number of common neighbours, here the number of recipes that contain
both. We change this projection in three ways. We keep only the ingredients
that occur in at least 50 recipes (362 of 695). We keep only the pairs that
share at least 10 recipes. We replace the number of shared recipes by the PMI
of the pair~\cite{church1990} and drop the pairs with PMI $\leq 0$. The result
has 362 nodes and 10{,}853 edges. It is a filtered, PMI-weighted variant of
the taught projection, not the taught projection itself, and PMI is an
extension that the course does not cover. The PMI of two ingredients $i$ and
$j$ is
\begin{equation*}
  \mathrm{PMI}(i,j) = \log_2 \frac{c_{ij}\, N}{c_i\, c_j},
\end{equation*}
where $c_{ij}$ is the number of recipes that contain both ingredients, $c_i$
and $c_j$ are the numbers of recipes that contain each of them, and
$N=45{,}749$ is the number of recipes with at least one ingredient. The
ratio compares how often the pair occurs with how often it would occur if the
two ingredients were independent. A positive PMI means that the pair is more
common than that; it is not a significance test. PMI tends to favour rare
pairs, and the threshold of 10 shared recipes limits this effect.

Shared-recipe counts tell us which pairs are common. PMI tells us which pairs
are more common than their frequencies predict. Table~\ref{tab:pmi} gives two
examples. Salt and garlic occur together in 7{,}571 recipes, but both are so
frequent that about 6{,}554 shared recipes are expected by chance, a ratio of
1.16 (PMI $0.21$). Ginger and soy sauce share fewer recipes (1{,}010) but
about 4.5 times more than expected (PMI $2.16$). On the same 10{,}853 edges,
Louvain finds 3 communities when the weights are the shared-recipe counts and
5 when they are the PMI values. We do not compare the modularity of the two
partitions, because the two weightings define different weighted graphs. We
report this comparison only to show that the weighting changes the groups.

\begin{table}[h]
\centering
\small
\begin{tabular}{@{}l r r r r r@{}}
\toprule
Pair & $c_i$ & $c_j$ & Together & Expected if independent & Ratio \\
\midrule
Salt -- Garlic & 18{,}745 & 15{,}995 & 7{,}571 & 6{,}554 & 1.16 \\
Ginger -- Soybean Sauce & 3{,}418 & 3{,}029 & 1{,}010 & 226 & 4.46 \\
\bottomrule
\end{tabular}
\caption{Recipes that contain both ingredients of a pair, compared with the
number expected if the two ingredients were independent, $c_i c_j / N$. The
ratio is the observed number divided by the expected number; its base-2
logarithm is the PMI. Two examples do not show that PMI is better for every
pair.}
\label{tab:pmi}
\end{table}
